% 简化对象图标的真实路径；调用方定义trust icon和trust wire。
\providecommand{\TrustVectorMail}[3]{%
 \begin{scope}[shift={(#1,#2)},scale=#3]
  \draw[trust icon,fill=white,rounded corners=.4mm] (0,0) rectangle (8,5.5);
  \draw[trust icon] (0,5.5)--(4,2.3)--(8,5.5) (0,0)--(2.7,2.4) (8,0)--(5.3,2.4);
 \end{scope}}
\providecommand{\TrustVectorFolder}[3]{%
 \begin{scope}[shift={(#1,#2)},scale=#3]
  \draw[trust icon,fill=FigureBlueFill,rounded corners=.4mm] (0,0)--(0,5)--(3,5)--(4,4)--(8,4)--(8,0)--cycle;
 \end{scope}}
\providecommand{\TrustVectorRecord}[3]{%
 \begin{scope}[shift={(#1,#2)},scale=#3]
  \draw[trust icon,fill=white] (0,0)--(6,0)--(6,6)--(4,8)--(0,8)--cycle;
  \draw[trust icon] (4,8)--(4,6)--(6,6);
  \draw[trust icon] (1,4.8)--(4.8,4.8) (1,3.1)--(4.8,3.1) (1,1.4)--(3.8,1.4);
 \end{scope}}
\providecommand{\TrustVectorDocs}[3]{%
 \begin{scope}[shift={(#1,#2)},scale=#3]
  \foreach \dx/\dy in {0/2,1.6/1,3.2/0}{
   \draw[trust icon,fill=white,rounded corners=.5mm] (\dx,\dy) rectangle ({\dx+6},{\dy+8});
  }
  \draw[trust icon] (4.2,6)--(8.2,6) (4.2,4.1)--(8.2,4.1) (4.2,2.2)--(7.2,2.2);
 \end{scope}}
\providecommand{\TrustVectorTrash}[3]{%
 \begin{scope}[shift={(#1,#2)},scale=#3]
  \draw[trust icon,fill=FigureRedFill] (.7,0)--(0,6)--(6,6)--(5.3,0)--cycle;
  \draw[trust icon,fill=white,rounded corners=.3mm] (-.4,6) rectangle (6.4,7);
  \draw[trust icon] (2,7)--(2,8)--(4,8)--(4,7) (1.8,1)--(1.4,5) (3,1)--(3,5) (4.2,1)--(4.6,5);
 \end{scope}}
\providecommand{\TrustVectorMesh}[3]{%
 \begin{scope}[shift={(#1,#2)},scale=#3]
  \foreach \xx/\yy/\id in {-5/-2/a,-5/2/b,0/-4/c,0/0/d,0/4/e,5/-2/f,5/2/g}{\coordinate (icon-\id) at (\xx,\yy);}
  \foreach \i in {a,b}{\foreach \j in {c,d,e}{\draw[trust wire] (icon-\i)--(icon-\j);}}
  \foreach \i in {c,d,e}{\foreach \j in {f,g}{\draw[trust wire] (icon-\i)--(icon-\j);}}
  % 层内同色、层间蓝红黄；连接保持灰色，不把整网染成单色。
  \foreach \id/\col in {a/FigureBlue,b/FigureBlue,c/FigureRed,d/FigureRed,e/FigureRed,f/FigureYellow,g/FigureYellow}{\draw[trust icon,fill=\col] (icon-\id) circle (.7);}
 \end{scope}}
